\chapter{消息传递与神经生物学}

“从根本上说，动物大致分为两类：动物，以及没有大脑的动物；后者被叫作植物。植物不需要神经系统，因为它们不会主动移动，它们不会把根拔起来，在森林大火中逃跑！任何能够主动移动的东西都需要神经系统；否则它很快就会死亡。”

——Rodolfo Llinas

\section{引言}

在第 4 章中，我们已经看到，对于两类生成模型，变分推断分别采取怎样的形式。本章聚焦于这些推断动力学所对应的过程理论：也就是解释大脑可能如何实现变分推断的理论。在这种贝叶斯信念更新的实现中，核心概念是贝叶斯消息传递，它涵盖了信念传播、变分消息传递以及其他相关方案。这些方案背后的基本思想是：并不是所有东西都直接依赖于所有其他东西。相反，生成模型中的每个变量只依赖于相对较少的其他变量。类似地，大脑也表现出稀疏的连接结构，其中任一神经元的活动只依赖于那些与它共享突触连接的神经元。本章关注的是，如何把与变分推断相关的稀疏消息传递，映射到生物计算中的稀疏连接结构上。

让我们暂时从第 4 章的技术细节后退一步，把注意力转向主动推断所伴随的过程理论。必须区分一个原理，也就是自由能最小化，与关于这一原理如何在某类系统中得到实现的过程理论，例如它如何在大脑中被实现。后者使我们能够提出可以接受经验数据检验的假设。为了讨论主动推断可能如何在大脑中显现，我们把第 4 章末尾看到的消息传递等同于突触通信，把梯度下降动力学等同于神经元活动。本章有两个目标：一方面，向具有技术背景的读者介绍主动推断的神经生物学；另一方面，向生物学家强调理论对于实践神经科学的意义。我们强调，本章并不打算给出关于主动推断过程理论的最终定论（Pezzulo, Rigoli, and Friston 2015, 2018; Friston and Buzsaki 2016; Friston and Herreros 2016; Friston, FitzGerald et al. 2017; Friston, Parr et al. 2017; Parr and Friston 2018b; Parr, Markovic et al. 2019）；它只是目前看来与现有证据最一致的一种解释。我们的目的也不是为某一个特定的过程理论背书，而是说明第 1--4 章提出的思想，如何被用于构造能够接受神经生物学测量检验的假设。

本章结构如下。首先，在第 5.2 节，我们讨论皮层微回路的作用。它是由若干彼此连接的神经元群体构成的功能单元，而这种连接模式在许多皮层区域中都会重复出现。我们将强调这种定型化回路与图 4.4 和图 4.6 中消息传递架构之间的关系，而后者本身也会在层级各层上反复出现。在第 5.3 节，我们转向效应系统，并讨论在主动推断框架下如何表述运动控制。这里处理的是运动皮层如何调节脊髓和脑干的反射弧，以生成有目的的行为。第 5.4 节将简要涉及丘脑、基底神经节等皮层下结构，它们在规划与决策中发挥重要作用。随后，在第 5.5 节，我们讨论突触效能的调制，包括神经递质在精度优化中的作用，以及学习中的可塑性变化。最后，在第 5.6 节，我们回到层级这一主题，并讨论决策与动作生成之间的关系。

\section{微回路与消息}

在第 4 章中，我们已经看到，介导变分推断的信念更新方程可以被解释为一种神经元网络。前面给出的连续模型与分类模型的推断方案，都会产生一种定型化回路，而这种结构会在层级式生成模型中重复出现。与此类似，大脑皮层的结构也具有一种定型化架构（Shipp 2007）。新皮层被划分为六层，或称板层，从表浅层（靠近脑表面）到深层（更接近皮层下白质）编号。每一层都以特定的细胞类型和连接模式为特征（Zeki and Shipp 1988; Felleman and Van Essen 1991; Callaway and Wiser 2009）；这些连接关系在图~5.1 的单个皮层柱示意图中得到了概括。

一个脑区中的皮层柱会与其他脑区的皮层柱以及皮层下结构发生连接。皮层区域通常被描绘为一个层级结构；粗略地说，最接近感觉输入或运动输出的区域位于层级的底部。随着我们远离这些区域，例如从初级视觉皮层走向次级视觉皮层，我们就在层级中向上移动。之所以能够提出这样的层级概念，是因为图~5.1 所示的层特异性连接结构。来自较低级皮层区域或感觉性（初级）丘脑核团的上行投射，往往指向第 IV 层的棘状星形细胞。较低级皮层区域的上行连接来源于其表浅锥体细胞，即第 II 和 III 层。来自较高级皮层区域的下行投射，会同时指向较低级区域的表浅层和深层，其来源则是深层，即第 VI 层的锥体细胞。此外，第 V 层中的深层锥体细胞（Betz 细胞）还会投射到其他靶点，包括基底神经节、次级丘脑核团等皮层下核团，以及脊髓运动神经元。

图~5.1 中间的示意图展示了一种可能的映射方式，即如何把预测编码网络（图 4.4）映射到皮层的层状解剖结构上（Friston, Parr, and de Vries 2017）。这一图有些复杂，但关键点如下：进入第 IV 层棘状星形细胞的上行输入，对应于隐藏原因的预测误差 $\epsilon_v^{(i)}$。来自第 III 层表浅锥体细胞的上行输出，则把同样的预测误差传递给层级中的上一层，即 $\epsilon_v^{(i+1)}$。下行输入表示来自更高层的预测 $g^{(i+1)}$，而下行输出表示对更低层的预测 $g^{(i)}$。在最低层，我们看到来自第 V 层的下行预测，这与左图中通向脊髓运动神经元的输出一致。我们将在第 5.3 节回到这一点。回忆第 4 章，隐藏原因的作用是把一个在多个不同时间尺度上运行的模型的各层级连接起来。这与隐藏状态不同，后者的作用是在某一特定时间尺度上形成动力学，这一点与图~5.1 中柱内的内在连接作用是一致的。

消息传递中的不对称性十分重要，因为它为上行活动与下行活动之间的差异提供了经验预测。其中一个预测是：我们可以预期这些消息由不同时间频率的神经活动来传递。原因在于，从期望计算预测误差所需的运算是非线性的（Friston 2019b）。这种非线性来自使用非线性函数 $g$ 来计算预测，而这通常会提高信号频率，例如对正弦波求平方会使频率翻倍。因此，一个可检验的预测是：来自误差单元的上行消息，可能出现在比来自期望单元的下行消息更高的频段中，如图~5.2 所示。这与实测到的频谱不对称性是一致的：上行连接通常与 $\gamma$ 频段相关，而下行连接则更常与 $\alpha$ 或 $\beta$ 频段相关（Arnal and Giraud 2012; Bastos, Litvak et al. 2015）。

图~5.1 右侧的示意图，则给出了把 POMDP 模型的消息传递解释为一种皮层微回路的方式。它与预测编码架构具有类似结构：期望 $s$ 由表浅和深层锥体细胞表示，并在皮层层级中向上和向下传播。此外，第 IV 层中的误差单元 $\epsilon$ 接收上行信号。不同于预测编码式架构的是，区域之间传递的消息是期望的混合，而不是误差（Friston, Rosch et al. 2017; Parr, Markovic et al. 2019）。不过，通过更新期望来最小化误差，也就是自由能梯度，这一总体结构仍然保持不变。这种消息传递区分了以策略为条件的期望（下标 $\pi$）与在各策略下平均后的期望。要从前者转换到后者，我们还需要关于策略的后验信念 $\pi$。我们会在第 5.4 节回到这一点；此处值得先指出的是，这种消息传递与某些皮层下结构投向皮层表浅层的事实是相容的，而这些皮层下结构可能正是用来计算这些信念的。

在图~5.1 中，要注意左图的解剖结构与中图、右图的消息传递方案之间并不存在一一对应关系。例如，中图与左图之间似乎存在一个差异：左图中进入该皮层柱的下行输入到达第 II 层和第 IV 层，而中图中对应的目标却是第 III 层。这表明，消息传递方案所蕴含的连接，并不一定体现为单一突触。指向第 III 层的下行抑制性连接，可能是通过如下组合来实现的：先有一条兴奋性投射到达第 IV 层的抑制性中间神经元，再由这些中间神经元投射到第 III 层。这个双突触通路就解决了两种架构之间看似存在的矛盾。

在第 5.3 节和第 5.4 节，我们将分别讨论第 V 层神经元在运动和规划中的作用，它们对应于通向脊髓（或脑干）与皮层下结构的投射。到了第 5.5 节和第 5.6 节，我们将讨论神经消息传递如何随时间被调制，并回到分类推断与连续推断这两类微回路之间的关系。

\section{运动指令}

图~5.1 左侧示意图表明，皮层第 V 层投射到脊髓锥体神经元，而这可以被解释为一种预测（Adams, Shipp, and Friston 2013）。图~5.3 将这一点进一步展开，展示了这一回路在脊髓中的组成部分。我们看到，一个基于初级运动皮层第 V 层 Betz 细胞所编码期望的预测，被送出并下传。该预测与传入脊髓背角的本体感觉输入相减，结果得到本体感觉预测误差。这一误差驱动肌肉活动，而肌肉活动又会导致这一误差自身被消除，因为本体感觉数据会改变，从而与预测相一致。把运动命令看作一种预测，是主动推断的核心思想之一，因为它凸显了行动与知觉之间的对偶性。行动这一部分，就是通过改变感觉数据来最小化预测与感觉数据之间的任何不一致。这意味着，要生成一个动作，我们唯一需要的，就是去预测：如果这个动作被执行，将会出现什么样的感觉后果。正因为本体感觉预测误差总是可以通过反射而非信念更新来解决，这也许解释了为什么初级运动皮层第 IV 层中颗粒细胞相对稀少（Shipp, Adams, and Friston 2013）。

这种运动控制的一个重要方面，是所谓的感觉衰减（Brown et al. 2013）。考虑这样一个问题：我们如何启动一个新动作？要做到这一点，我们必须能够预测“自己正在移动”。然而，在我们真正开始移动之前，手头的本体感觉数据其实与这一假设相矛盾，因此可能促使我们修正它。因此，我们需要一种机制，能够阻止感觉数据去更新我们的期望，好让我们在最初事实上仍然为假的情况下，也能保持“我正在移动”这一信念，并最终通过行动把它实现出来，这与观念运动现象相呼应。其含义是，我们必须能够通过降低本体感觉数据的增益，把注意力从这些数据上移开。从技术上说，这个增益由这些数据被预测时所对应的精度，也就是方差的倒数，来给定。要实现这种衰减，我们就需要下行运动通路不仅预测数据本身，还要预测这些数据的精度，或者说对这些数据赋予的信心，并在启动动作时降低这种精度。这就是所谓的感觉衰减。它可以被看作感觉注意的补足面，使我们有能力忽略某些预测误差，例如由扫视眼动产生的误差；在这种情形下，感觉衰减被称为扫视抑制。无法正常衰减被认为是多种神经系统和精神病理综合征的核心机制，包括被动体验现象（Pareés et al. 2014），以及更极端的精神分裂症紧张状态和帕金森病中的动作启动失败。

\section{皮层下结构}

除了投向脊髓之外，皮层第 V 层还会投射到若干其他结构，其中之一便是纹状体（Shipp 2007; Wall et al. 2013）。纹状体位于大脑深部，由尾状核和壳核构成，是一组复杂网络即基底神经节的输入核团。中型棘状神经元是纹状体的功能单元，它们接收来自皮层的输入，并投射到基底神经节的其他核团。这些神经元分为两类：表达 D1 多巴胺受体的，以及表达 D2 受体的。多巴胺会增强前者的活动，并减弱后者的活动。前者构成了基底神经节中的直接通路，通过一次抑制性突触与输出核团，即苍白球内侧部和黑质网状部，相连接。D2 细胞则形成间接通路，其路径稍微更复杂，包含两次抑制性突触和一次兴奋性突触。纹状体抑制苍白球外侧部，而后者又抑制丘脑底核（STN）。丘脑底核再投射到基底神经节的输出核团。因此，输出核团受到直接通路的抑制，而受到间接通路的去抑制。由于这些输出核团本身也是抑制性的，所以激活表达 D1 的纹状体神经元，净效果就是解除对行为的抑制；而表达 D2 的神经元则会压制这种行为（Freeze et al. 2013）。

既然我们已经把投向脊髓的投射解释为本体感觉预测，那么皮层第 V 层投向纹状体的投射，又应该对应什么样的消息呢？观察图~5.1 可以得到一种可能的答案。在这一层中，预测结果 $o$ 与偏好结果和预测结果之间的差异 $\varsigma$ 一起出现，二者结合起来计算某一策略的期望自由能 $G$。把这一计算放在纹状体中，是与“基底神经节参与规划”这一观念相一致的，也就是它评估不同的行动方案。图~5.4 展示了一种可能的映射方式：如何把策略评估的消息传递映射到基底神经节的解剖结构上。

从图~5.4 中应当抓住的关键点是，正如第 4 章所述，关于策略的后验概率 $\pi$，在图中显示于苍白球内侧部，是根据这些策略的期望自由能来计算的。这一模式沿着直接通路传播：从皮层第 V 层出发，经由纹状体，到达基底神经节输出核团。不过，其中还有一些额外的细节。在图 4.4 左侧的层级架构中，我们看到，高层状态上的期望能够影响低层策略上的信念。图~5.4 左侧对此作了展示：高层状态下的期望观测，被用来形成经验先验 $E$，这些先验会独立于期望自由能而影响策略选择。我们将在第 7 章回到这一点，但核心思想是：在特定情境下，我们会对自己如何行动持有某种信念。当我们再次处于同样情境时，这些先验期望往往会对策略评估形成偏置，这很像习惯的形成。从这个意义上说，$E$ 与 $G$ 的影响可以分别被理解为习惯驱动与目标导向驱动。在强化学习中（Lee et al. 2014），它们有时分别被称为“无模型”与“有模型”系统。\footnote{这里的“无模型”和“有模型”并不是说前者完全没有内部模型，而是指其行为更依赖经验累积的先验偏置，后者则更依赖对结果与路径的显式评估。} 如果把它们分别与基底神经节的直接通路和间接通路对应起来，就会产生一个有趣后果：这意味着，多巴胺会调节这两种驱动之间的平衡。请记住，多巴胺倾向于促进直接通路，从而推动特定策略的执行（Moss and Bolam 2008），而这些策略大概正是那些具有最低期望自由能的策略。相反，在低多巴胺状态下，我们也许更应预期由间接通路承载的、对情境敏感的先验会占优势，而它们的作用正是在特定情境中压制那些不可信的策略。从某种意义上说，纹状体中的多巴胺可以被看成是在调节“推断该做什么”与“推断不该做什么”之间的平衡（Parr 2020）。

上述解释与多巴胺系统受到扰动时的表现是一致的。严重帕金森病中的多巴胺耗竭，会导致运动不能，也就是无法执行特定策略；而外源性多巴胺激动剂则会促进冲动性行为（Frank 2005; Galea et al. 2012; Friston, Schwartenbeck et al. 2014）。此外，这也与基底神经节功能的概念模型相一致。例如，Nambu（2004）提出，直接通路会造成对苍白球内侧部快速而聚焦的抑制，随后出现广泛而缓慢的兴奋，这分别导致基底神经节靶结构的兴奋和抑制。这被认为能够形成一种“中心-周边”式模式，从而促进高特异性的运动程序。这与如下图景是一致的：快速过程通过计算期望自由能来促进动作，而较慢通路提供的更广泛情境化信息，则用于传递经验先验。

关于图~5.4 最后还需要指出的是：同一个基底神经节回路，接收到了来自皮层层级两个层次的输入，上标对此作了区分。这暗示了这样一种情况：在指向间接通路神经元时，来自时间上较慢区域的影响占优；而直接通路则同时受到快与慢两种影响。随着我们在皮层层级中向上移动，神经元通常会表征更慢的动力学。例如，我们预期额叶皮层区域表征的时间尺度会长于顶叶区域。这与皮层输入在基底神经节各通路上的解剖分布相一致（Wall et al. 2013）。这种在间接通路中的时间粗粒化，还伴随着空间上的粗粒化：直接通路的中型棘状神经元具有更大的树突分支范围（Gertler et al. 2008），从而能够实现更精细的调节。因此，图~5.4 所示的解剖构造，同时得到了临床病理学证据，例如帕金森综合征，以及细胞形态学证据的支持。

除了基底神经节之外，还有许多其他重要的皮层下结构也参与神经消息传递。下一节我们将讨论神经调质系统的起源结构；在本节末尾，我们先简要谈及丘脑。我们无法在这里对这一高度复杂的结构做出充分讨论，但可以概述一些基本原则。丘脑通常被分为初级核团和次级核团。图~5.1 表明，初级丘脑核团可以发挥与低级皮层区域相同的作用，也就是说，它们投向皮层第 IV 层，并接收来自第 VI 层深层锥体细胞的输入（Thomson 2010; Olsen et al. 2012）。例如，视觉系统中的外侧膝状体，就常被看作眼睛与视觉皮层之间的中继站。与那些表征预测误差的神经元类似，它既接收来自眼睛的感觉信息，也接收来自皮层的反馈投射，而后者可以被理解为预测。二级丘脑核团则包括内侧背核和枕核，它们分别与额叶皮层和后部皮层发生相互作用。这些结构可能参与对感觉输入或更高阶输入的二阶统计量，也就是精度和方差，进行预测，并且被发现与图-底分离任务有关（Kanai et al. 2015）。如果做一个简化理解，那么丘脑划分为初级核团和次级核团，可能正对应于一阶统计量与二阶统计量之间的划分。

\section{神经调制与学习}

结构性神经解剖很重要，但它只能告诉我们神经处理图景的一部分，因为一条连接的存在，并不能充分说明它究竟是如何被使用的。举例来说，考虑图~5.4 中黑质的作用。它对纹状体连接施加的调制效应，会根据释放的多巴胺量不同，使基底神经节产生非常不同的输出。像这样的突触效能快速调制，可以与学习中发生的、更慢但更持久的变化形成对照。本节聚焦于突触效能发生改变的这两种方式。

精度是理解神经调制的一个重要概念（Feldman and Friston 2010）。我们在第 4 章以及前面关于感觉衰减的讨论中已经接触过这个概念；在那里我们看到，精度作为预测误差上的一个乘性权重而发挥作用。更广泛地说，精度是一个概率分布中“我们有多大把握”的度量。二者之间的关系很简单。如果我们对于“数据如何由隐藏状态生成”持有非常高精度的信念，那么与在我们对这些信念缺乏把握时相比，观察到数据将更强烈地更新我们关于隐藏状态的信念。当这种信念更新表现为神经元放电变化时，一个更高精度的似然分布，就体现为对给定感觉刺激更强的神经反应。这一点对于从注意（Parr and Friston 2019a）到多感觉整合（Limanowski and Friston 2019）的一系列认知功能都至关重要。

这种关于突触增益控制的视角，告诉我们一件简单却重要的事：如果精度是生成模型中某个分布的属性，那么不同分布就应当对应不同的精度。这在直觉上是合理的，因为我们对于感觉的可靠性、事物将如何随时间演化，甚至我们可能如何行动，都可能有高低不等的信心（Parr and Friston 2017b）。在基底神经节的情境中，我们已经看到其中最后一种，即关于策略的精度。与之相关的精度越高，我们就越有信心认为自己的策略会最小化期望自由能。

如果多巴胺系统，也就是起源于中脑黑质致密部和腹侧被盖区的系统，其功能之一是传递“该做什么”的信心，那么其他神经调质系统是否也扮演类似角色？表~5.1 总结了把不同精度与不同神经调质系统联系起来的证据，并给出了这些精度有时使用的符号。具体而言，起源于 Meynert 基底核的胆碱能系统，似乎传递某些似然分布的精度。

\paragraph{表 5.1：主动推断中神经递质的可能作用}

乙酰胆碱对应似然精度 $\zeta$。相关证据包括：丘脑-皮层传入纤维上存在突触前受体（Sahin et al. 1992; Lavine et al. 1997）；它能够调节视觉诱发反应的增益（Gil et al. 1997; Disney et al. 2007）；药理学操纵会改变有效连接（Moran et al. 2013）；在药理操纵条件下，对行为反应的建模结果也支持这一点（Vossel et al. 2014; Marshall et al. 2016）。

去甲肾上腺素对应转移精度 $\omega$。证据包括：前额叶持续性活动，也就是延迟期活动，而这种活动要求转移概率具有足够高的精度，其维持依赖于去甲肾上腺素（Arnsten and Li 2005; Zhang et al. 2013）；瞳孔对令人惊讶，也就是低精度序列的反应与其相关（Nassar et al. 2012; Lavín et al. 2013; Liao et al. 2016; Krishnamurthy et al. 2017; Vincent et al. 2019）；药理操纵条件下的行为建模同样提供支持（Marshall et al. 2016）。

多巴胺对应策略精度 $\gamma$。证据包括：其受体在纹状体中型棘状神经元上以后突触形式表达（Freund et al. 1984; Yager et al. 2015）；计算性 fMRI 揭示，当精度发生变化时，中脑活动也随之变化（Schwartenbeck, FitzGerald, Mathys, Dolan, and Friston 2015）；药理操纵条件下的行为建模也支持这种联系（Marshall et al. 2016）。

血清素则可能对应偏好精度，或者内感受似然精度 $\chi$。证据包括：其受体在内侧前额叶皮层第 V 层锥体细胞上表达（Aghajanian and Marek 1999; Lambe et al. 2000; Elliott et al. 2018）；而内侧前额叶皮层区域又与内感受加工和自主神经调节高度相关（Marek et al. 2013; Mukherjee et al. 2016）。

上述内容据 Parr and Friston 2018 整理。

来自蓝斑的去甲肾上腺素系统，似乎在传递时间转移的精度方面发挥作用。血清素系统的情况则没有那么明确，但它可能与先验偏好的精度有关。

为什么把这些精度与神经调质系统对应起来会有用？答案有三点：它能帮助我们解释已观察到的生物学现象，提出新的假设，并发展无创方式来测量精度。下面我们分别举一个例子。首先，在解释已观察生物学方面，关于多巴胺的经验测量，经典上看起来像是“奖励预测误差”（Schultz 1997）：动物在获得意外果汁奖励，或看到预示果汁即将出现的线索时，多巴胺会升高。主动推断对这些发现提供了一种替代性解释（Schwartenbeck, FitzGerald, Mathys, Dolan, and Friston 2015）。当我们获得奖励，或满足自身偏好，又或者遇到预示未来奖励的线索时，我们对自己正在追随一项能够最小化期望自由能的策略这件事，会更有信心。这种信心的提升，就表现为多巴胺的尖峰。

其次，在提出假设方面，一个例子是路易体痴呆中胆碱能信号的下降（Parr, Benrimoh et al. 2018），这种疾病会导致复杂视觉幻觉。一种合理解释是：较高阶视觉皮层中的病理积累，促使这些区域发出的预测与初级视觉皮层中的活动之间出现不匹配。这种不匹配会降低相关似然分布的置信度，从而导致胆碱能信号丧失。精度丧失的后果，是个体无法根据感觉数据来更新信念，也就是知觉失去了感觉所提供的约束。这或许可以解释为什么这一疾病会发展出幻觉性知觉。

第三，在无创测量精度参数方面，一个例子是计算型表型的识别。中枢神经化学活动在外周有许多表现形式，例如自发眨眼率与多巴胺之间的关系（Karson 1983），以及瞳孔大小与去甲肾上腺素之间的关系（Koss 1986）。近期关于后者的研究（Vincent et al. 2019）已经表明，一名理想贝叶斯观察者本应推断出的转移精度，与瞳孔收缩和扩张的动力学之间存在对应关系。这意味着，我们或许可以通过此类外周测量，去探测某个人隐含生成模型中的经验先验信念。

虽然精度的快速变化很重要，但它是一种相当粗糙的有效连接优化方式；它只会提高或降低信号增益，而无法实现更细腻的调整。如果我们想改变信号被解释的方式，就必须依赖学习。关于这一点，我们将在第 7 章中详细讨论。这里的基本思想是：我们不仅持有关于世界状态的信念，也持有关于固定参数，或者缓慢变化参数的信念，而这些参数决定了变量之间的依赖关系（Friston, FitzGerald et al. 2016）。这些信念的物质基础，是表征时间变化变量，如隐藏状态或结果的神经群体之间的突触连接效能。当我们观察到一个结果，并相信它是由某个给定状态生成时，我们就能够更新连接二者的参数信念，这体现为它们未来共同出现概率的上升。换句话说，连接这两组神经元的突触会得到加强。这就得到了 Hebb 那句著名命题的含义，可概括为：“一起放电的细胞，会连接在一起。”

图~5.1 的一个重要特征在于：无论是在预测编码方案还是边缘消息传递方案中，进入和离开皮层柱的连接都与似然分布有关。相比之下，转移概率和连续动力学则依赖于微回路内部的连接。这暗示着：学习动力学应当导致内在连接发生变化，而学习观测模型则应当修改外在连接。利用动态因果建模等技术，我们可以从神经影像数据中评估有效连接，并据此检验这些假设（Tsvetanov et al. 2016; Zhou et al. 2018）。这正凸显了此类过程理论的作用：它让我们能够超越抽象理论化，形成具体而可检验的假设。

\section{连续层级与离散层级}

最后，值得强调的是，在神经层级的低层中，表征往往是连续的，而在更高层中，则会转向分类变量。原因在于，我们前面讨论过的离散消息传递方案与连续消息传递方案，很可能在大脑中同时存在，因为我们不仅能够持有分类式信念，例如识别一个物体是什么、一个人是谁，也能够与连续变化的感觉受体和效应器接口，例如肌肉长度或视觉亮度对比。这一点也反映在神经生理学中：有些神经元对特定刺激具有选择性，而另一些神经元的反应则会随着刺激强度成比例变化。

一个有趣的观察是，我们与周围世界的接口处于连续域中，而这意味着，大脑中任何层级的最低层都必须是连续的。尽管如此，我们在图~5.4 中看到，基底神经节中的策略选择可以被表述为一个离散过程，也就是在不同动作之间进行选择。这说明，我们可以把动作理解为若干离散轨迹被组合成有目的行为。最低层也许处理的是肌肉长度所需发生的具体变化，而下行输入则建立在“该做哪一种动作”的决策之上。从生成模型的角度看，这意味着：我们要把关于世界的不同离散假设，与这些假设所蕴含的连续动力学联系起来。第 8 章中我们会回到这个问题，讨论如何从计算角度把这二者整合起来。在这里，我们只需指出：离感觉受体越远，神经系统中就越容易出现离散化表征。事实上，神经生理学中经典感受野的存在，本身就可以被理解为一种概率性表征：它表示世界正处于知觉状态空间中的某个特定区间，而这个状态空间被感受野所铺砌，因此被切分为大量小类别。图~5.5 将这些方案汇集到一起，可视作本章思想的总结。

\section{总结}

本章试图勾勒出第 4 章生成模型所蕴含的消息传递方案，与推断、行动和规划的神经生物学之间的连接点。把消息传递与神经元通信联系起来，我们能得到什么？它使我们能够根据我们假定大脑正在反演的生成模型，提出经验性预测。这种预测可以表现为诱发反应，也就是当大脑接收到感觉刺激时，头皮上可测得的电位变化，而其时间进程将取决于该刺激诱发了多大程度的信念更新。另一种路径，则是使用计算神经影像方法，把模拟得到的推断过程，与那些表现出相似时间动力学的脑区联系起来（Schwartenbeck, FitzGerald, Mathys, Dolan, and Friston 2015）。建立这种联系，对于理解计算性，也就是神经学和精神病学意义上的病理，以及相应治疗，都很重要，因为它使我们能够用生物学语言来表达功能性病理。

最后，值得承认的是，本章明显没有涉及大脑的很多部分，这部分是因为篇幅限制，另一部分则是因为神经科学本身仍在发展中。对于本章提出的解释框架，还存在许多可以扩展，甚至替代它的机会。在一定程度上，我们可以从已经看到的内容向外推演。例如，杏仁核的某些部分在细胞构筑学上与基底神经节核团是等价的。这是否意味着，杏仁核也在评估某一类策略？对于自主神经策略来说，它是否扮演着类似基底神经节在骨骼运动领域中的角色？其他结构，例如枕核，是否会对其他类型，例如心理层面的策略，发挥类似作用？对于那些不同于图~5.1 六层结构的皮层架构，我们又该如何理解？小脑和海马结构都呈现出各自独特但又具有定型性的微回路（Wesson and Wilson 2011）。我们应当把它们看成是同一套贝叶斯消息传递方案在解剖上的重新排列，还是说它们处理的是生成模型中的其他方面（Pezzulo, Kemere, and van der Meer 2017; Stoianov et al. 2020）？提出这些问题，并不是为了在这里给出答案，而是为了指出理论神经生物学中一些令人兴奋的未来研究方向。主动推断及其相关过程理论，为处理这些问题提供了一个严格的形式化框架和概念框架。
